% SPDX-License-Identifier: GPL-3.0-or-later
% Hindi (हिन्दी) -- what the reading editions' preamble leaves to the language.
% Input by lib/frank-preamble.tex as \input{\FrankLib/lang/\BookLang.tex};
% docs/languages.md section 6 lists the hooks every file here must define and
% section 12 the ones it may; the guide's page "Adding a language"
% (html-guide/markdown/lookup-and-languages/adding-a-language.md) walks
% through writing one.
%
% Scaffolded from lib/lang/_template.tex by lib/newlang.py.  Everything below
% the switches came out of lib/languages.json, so it agrees with the registry
% by construction; \FrankHowTo did not, and is the one thing you must write.
% The model to read while writing it is lib/lang/fa.tex.

% ---- the babel language and the switches ----------------------------------
% \FrankLangName is what \foreignlanguage and \begin{otherlanguage} are given
% everywhere in the core preamble: babel must know this name, or every passage
% comes out in the roman with a "Unknown language" error.
\newcommand{\FrankLangName}{hindi}
\FrankRTLfalse           % chunk left, gloss right; \pw needs no bidi isolate
\FrankHasBarefalse       % nothing comes off, so pass 3 would repeat pass 1
\FrankHasAltfalse        % whether \FrankAltFont exists (the title, the chapter number)
\FrankHasReadingfalse    % whether \chr's third argument is set as a reading
\FrankHasVertfalse       % whether pass 4 is \FrankVertPass rather than the font pass

% ---- fonts ------------------------------------------------------------------
% No face is NAMED here: \FrankPickFont walks the registry's tex.main and then
% tex.main_fallbacks and takes the first one installed, so the .tex and the
% .json cannot drift apart (docs/languages.md section 12).  A bundled face --
% one that travels in lib/fonts/ -- is found because build.sh puts that
% directory on OSFONTDIR.
\babelprovide[import=hi,onchar=ids fonts]{hindi}
\FrankPickFont{\FrankMainFontName}{hi}{main}{main_fallbacks}
\babelfont[hindi]{rm}[Script=Devanagari,Language=Hindi,Renderer=HarfBuzz]{\FrankMainFontName}

% ---- the vocabulary line -----------------------------------------------------
% \vb{infinitive}{sound}{stem}{sound}{perfective}{sound}{meaning}: these two are
% the labels printed before the second and third forms.  The stem is the
% infinitive less -nā; the perfective, masculine singular, is why the entry
% exists -- करना किया, जाना गया, होना हुआ, देना दिया, लेना लिया -- and a regular one
% is given too, because a reader who has met three irregulars does not yet
% know which verbs are regular.  docs/lang/hi.md names the same three forms in
% the same order, or the annotator fills slots the page labels differently.
% What the three do not say is whether the verb's subject takes ने in the
% perfective -- the one thing a reader needs to build a past sentence, and
% the one the forms hide.  It goes in ONE parenthesis after the meaning:
%   (+\pw{ने})     it takes ने: \vb{करना}{karnā}{कर}{kar}{किया}{kiyā}{to do (+\pw{ने})}
%   (±\pw{ने})     it goes either way (समझना, बदलना)
%   nothing       it never does (जाना, आना, हँसना)
% The mark is decided from a small hand-kept table first (lib/lang/hi.verbs.json)
% and from the dictionary's transitivity only after it: on 152 common verbs the
% dictionary's code agreed with ने for 110, contradicted it for 9 and was
% missing for 33 (करना, जाना and आना among them), which is too many to print
% unread.  A conjunct's light verb keeps its mark with no meaning before it:
% \vb{करना}{karnā}{कर}{kar}{किया}{kiyā}{(+\pw{ने})}\bw{काम}{kām}{m. work}.
% The pair below comes from the registry's "vb_labels" ("stem", "perf."), which is
% where the reader takes it from as well, and the three forms in words are its
% "vb_forms": change one and change the other, or the PDF and the reader of
% one book label the same form differently.  lib/newlang.py --check compares them.
\newcommand{\FrankVbPres}{stem}
\newcommand{\FrankVbPast}{perf.}

% ---- the two Lua filters ---------------------------------------------------
% frank_strip(s): the bare form of a chunk -- what pass 3 shows.  Written from
% the registry's `strip`; the identity when there is nothing to take off.
% frank_latin(s): the label's digits as ASCII, because a PDF destination name
% cannot carry the language's own figures.  Written from `digits`.
% Both are called for every language, so both must exist even when they do
% nothing.  No #, % or literal backslash in a \directlua body: TeX reads it
% first (NOTES section 9 of the Persian book).
\directlua{
  function frank_strip(s) return s end
  function frank_latin(s)
    local t, n = {}, 0
    for _, c in utf8.codes(s) do
      if c >= 0x0966 and c <= 0x096F then c = c - 0x0966 + 0x30 end
      n = n + 1 ; t[n] = utf8.char(c)
    end
    return table.concat(t)
  end
}

% ---- the "How to read this" page -----------------------------------------------
% Printed by lib/frank-frontmatter.tex under its heading.  \pw is defined by the
% core preamble, after this file is read; it is only expanded when the page is set.
\newcommand{\FrankHowTo}{%
Every passage comes twice.

\smallskip
\textbf{First} the whole sentence in Devanagari, and nothing else. Try to read
it. Get as far as you can and do not worry about what you cannot get — the
point is the attempt, made before any help arrives.

\smallskip
\textbf{Second} the same sentence in small chunks, Hindi on one side and on
the other three lines: how it sounds, then how to look each word up, then what
it means. Now you find out how much you had right. There is no third pass:
Devanagari writes its vowels, so the sentence you attempted is already the
sentence as it is printed, with nothing held back.

\smallskip
The line of sound is a \emph{pronunciation} and not a spelling. Devanagari
writes a vowel \textit{a} inside every consonant, and Hindi very often does
not say it: \pw{कमल} is \textit{kamal} and not \textit{kamala},
\pw{समझना} is \textit{samajhn\=a}. Trust the line over the letters — it is
where the word's real shape is.

\smallskip
\textit{\=a \=\i{} \=u} are long, held about twice as long as the short ones,
and the difference carries meaning. A dot under a letter means the tongue
curls back to the roof of the mouth: \textit{\d{t} \d{d} \d{n} \d{r}} are the
retroflex sounds, and the plain \textit{t d n} are made with the tongue on the
teeth, further forward than the English ones. An \textit{h} after a consonant
is a puff of breath and not a separate sound — \textit{kh gh ch jh th dh ph
bh} are single letters. \textit{\'{s}} is \textit{sh}, \textit{c} is the
\textit{ch} of \textit{church}, \textit{\.{n}} and \textit{\~{n}} the nasals
of \textit{sing} and \textit{onion}, and \textit{\.{m}} says the vowel before
it is spoken through the nose. \textit{q x \.{g} z f} appear in words that came
from Persian and Arabic, and are printed only where the text itself marks them
with a dot.

\smallskip
Verbs are given as infinitive, stem, perfective, meaning — the perfective is
the one you cannot guess (\pw{करना} \textit{karn\=a}, \pw{किया}
\textit{kiy\=a}; \pw{जाना} \textit{j\=an\=a}, \pw{गया} \textit{gay\=a}). A
\textit{+}\pw{ने} in brackets after the meaning says that in the perfective the
subject takes \pw{ने} and the verb agrees with the object instead
(\pw{लड़के ने किताब पढ़ी}, the boy read the book: \pw{पढ़ी} is feminine, for the
book); \textit{±}\pw{ने} that it goes either way; nothing, that it never does.
Every noun carries \textit{m.} or \textit{f.}: the gender is invisible in the
word and everything in the sentence agrees with it.
}
